% Generated draft: requires mathematical and source review.
\documentclass[11pt]{article}
\usepackage[margin=25mm]{geometry}
\usepackage{fontspec}
\setmainfont{Latin Modern Roman}
\usepackage{amsmath,amssymb,mathtools,unicode-math}
\setmathfont{Latin Modern Math}
\usepackage{enumitem,needspace,xurl,hyperref}
\setlength{\parindent}{0pt}
\setlength{\parskip}{5pt}
\setlength{\emergencystretch}{3em}
\newcommand{\N}{\mathbb{N}}
\newcommand{\Z}{\mathbb{Z}}
\newcommand{\Q}{\mathbb{Q}}
\newcommand{\R}{\mathbb{R}}
\newcommand{\C}{\mathbb{C}}
\newcommand{\F}{\mathbb{F}}
\newcommand{\A}{\mathbb{A}}
\newcommand{\PP}{\mathbb{P}}
\newcommand{\E}{\mathbb{E}}
\newcommand{\eps}{\varepsilon}
\newcommand{\dd}{\,\mathrm{d}}
\begin{document}
\section*{UnsolvedMath ID: 30003252}
\textbf{Source problem code:} catalogue entry
\subsection{Definitions and mathematical statement}
Let \(S=\{n^2:n\in\Z\}\) be the set of integer squares. For a positive integer \(M\), a sequence \((a_1,\ldots,a_M)\in S^M\) has constant second difference \(2\) if
\[
a_{i+2}-2a_{i+1}+a_i=2
\qquad\text{for every integer }i\text{ with }1\le i\le M-2.
\]
Call such a sequence consecutive-square in form if there is an integer \(x\) such that \(a_i=(x+i)^2\) for every \(1\le i\le M\). This concerns the square values; no choice of signs for their square roots is prescribed.

The supplied formulation of Büchi's problem asks whether
\[
\exists M\in\Z_{\ge1}\;\forall(a_1,\ldots,a_M)\in S^M,\qquad
\left[
\forall i\in\Z,\quad
1\le i\le M-2\ \Longrightarrow\ a_{i+2}-2a_{i+1}+a_i=2
\right]
\Longrightarrow
\left[
\exists x\in\Z\;\forall i\in\{1,\ldots,M\},\quad a_i=(x+i)^2
\right].
\]
Here \(M\) must be uniform over all sequences of that length, whereas \(x\) may depend on the sequence.

\emph{Unresolved qualification.} The supplied statement allows every positive length \(M\), including lengths below \(3\), for which the second-difference condition has an empty index range. This leaves a substantive short-length issue in the stated target. A lower bound ensuring a nonempty second-difference condition is not supplied; none has been silently added. The intended nonvacuous formulation requires human mathematical review.
\subsection{Short English statement}
The supplied question asks for a fixed length \(M\) such that every sequence of \(M\) integer squares with second differences \(2\) consists of squares of consecutive integers. The allowed short lengths leave an unresolved formulation issue. The supplied partially solved label and dated literature assessment are curation claims, not independently verified conclusions.
\subsection{Sources}
\begin{itemize}
\item Local UnsolvedMath registry, supplied record with permanent ID \(30003252\), titled ``Büchi's Problem for Sequences of Integer Squares,'' source problem code \texttt{catalogue entry}, in the Oberwolfach Reports Open Problems collection. No individual proposer is named; the supplied proposed year is \(2016\).
\item Supplied report citation: \emph{Definability and Decidability Problems in Number Theory} (2017), Oberwolfach Reports, pp.~2793--2866, DOI \texttt{10.4171/owr/2016/49}. \url{https://doi.org/10.4171/owr/2016/49}. The supplied background describes these reports as records of mathematical workshops and their open-problem discussions.
\item Supplied literature reference: ``Recent progress on Büchi's problem.'' \url{https://arxiv.org/abs/2412.16740}.
\item Supplied literature reference: ``Pasten--Pheidas--Vidaux survey.'' \url{https://www.mathnet.ru/rus/znsl3818}.
\item Literature assessment dated exactly \texttt{2026-08-21}: the supplied curation labels the problem partially solved and states that Büchi's problem remains open in general, with recent nontrivial finite-length results including a five-square result. This is supplied dated curation, not independently verified. Its source note explicitly requests independent verification of time-sensitive status and cited claims; the precise hypotheses of the five-square result are not included in the record.
\end{itemize}
\end{document}
